\chapter{[MP1] Medizinprodukte 1}
\label{chp:GER}
\gAasRsN{RS~IX}
\emph{Maintainer: Ing. Roman Koch}

\minitoc

Das folgende Kapitel soll einen Überblick über die wichtigsten Funktionen der
entsprechnden Geräte biten. Es ersetzt keine vom Hersteller herausgegebene
Gebrauchsanweisung. Diese ist vom Betreiber dem Personal in geeigneter Form
zugänglich zu machen und soll von den Anwendern konsultiert werden.

\section{Medizinproduktegesetz - MPG}

\gMarried{
Das Medizinproduktegesetz regelt die Funktionstüchtigkeit,
Sicherheit, Inbetriebnahme, Instandhaltung, den Betrieb und
die Anwendung, Überwachung, Sterilisation, Desinfektion und
Reinigung von Medizinprodukten. Als Medizinprodukte gelten
alle Instrumente, Apparate, Vorrichtungen, Stoffe oder
andere Gegenstände (inkl. Software) die vom Hersteller zur
Anwendung am Menschen bestimmt sind. Es gibt über 400.000
verschiedene Produkte. Doch eines haben alle gemeinsam: die
\gE{CE-Kennzeichnung}.
}{
\centering \includegraphics[width=0.5\linewidth]{./images/geraete/Ce-logo.jpg}
\gCaptionof{figure}{Das CE-Kennzeichen.}
}


\paragraph{Einteilung in Risikoklassen}
\begin{table}[H]
\gCaptionof{table}{MPG-Risikoklassen.}

\begin{tabulary}{\linewidth}{lLL}\addlinespace
\gTabSub{Klasse I}
&
Kein oder geringes Anwendungsrisiko
&
z.B: Verbandsmaterial, Gehhilfe, Rollstuhl, Tragsessel, Sehhilfe, Stethoskop, Trage, Beatmungsbeutel, Fieberthermometer, Druckminderer
\\\addlinespace
\gTabSub{Klasse IIa}
&
Anwendungsrisiko ohne entsprechender Ausbildung bzw. Unterweisung
&
z.B: aktive Diagnosegeräte, Absaugeinheit und -katheter, EKG-Kabel, Zuckerstreifen, Trachealtuben, Dauerkatheter
\\\addlinespace
\gTabSub{Klasse IIb}
&
Erhöhtes Anwendungsrisiko ohne entsprechender Ausbildung bzw. Unterweisung
&
z.B: Beatmungsgerät, Defibrillator, Pulsoxymeter, Babyinkubator, Röntgengeräte
\\\addlinespace
\gTabSub{Klasse III}
&
Besonders hohes Anwendungsrisiko ohne entsprechende Ausbildung bzw. Unterweisung
&
z.B: Perfusor, Insulinpumpe, Herzkatheter oder -klappen, Medizinprodukte mit Arzneimittelkomponenten
\\\addlinespace\bottomrule
\end{tabulary}
\end{table}

\paragraph{Pflichten des Sanitäters gemäß MPG}

Sanitäter sind lt. \gAbk{MPG} verpflichtet, die
Medizinprodukte richtig und sicher einzusetzen bzw. anzuwenden, so wie es
ihnen bei ihrer Ausbildung bzw. Unterweisung gelehrt wurde.
Außerdem sind sie verpflichtet die Geräte vor dem Gebrauch
mittels eines Gerätecheck auf ihre Funktion zu überprüfen.

% In diesem Abschnitt werden weitere Maßnahmen und Geräte für die Diagnostik und
% Patientenüberwachung (Monitoring) besprochen. Manche Maßnahmen werden in anderen
% Abschnitten besprochen (z.B. Blutdruckmessung
% (\prettyref{topic:blutdruckmessung}), Messung der Herzfrequenz (Pulsmessung,
% \prettyref{topic:pulsmessung}), \ldots). 

\section{Diagnostik und Monitoring}


\subsection{Stethoskop}
% TODO : Neufassung
\begin{anfxerror}[author=GABS]{Fehlender Text.}
Text fehlt! Hier muß neuer Text geschrieben werden!
\end{anfxerror}

\subsection{Manuelle Blutdruckmessgeräte}
% TODO : Neufassung
\begin{anfxerror}[author=GABS]{Fehlender Text.}
Text fehlt! Hier muß neuer Text geschrieben werden!
\end{anfxerror}




\subsection{EKG}
\label{topic:ekg}
\index{EKG}

Das Elektrokardiogramm (EKG) zeigt mittels eines entsprechenden
Gerätes die elektrische Herzaktivität des
Reizbildungs- und Reizleitungssystems an. Dazu werden auf
dem Körper nach einem bestimmten Muster Elektroden
aufgeklebt und mittels eines Kabels mit dem EKG-Gerät
verbunden. Im Rettungsdienst werden fast ausschließlich
Klebeelektroden für den einmaligen Gebrauch verwendet. 

\gAuthNotes{GABS: Sollte im Defi-Teil durchgenommen
werden?:} \gZ{EKG -- Ablauf einer Erregung

\begin{itemize}
    \item P-Welle $\approx$ Vorhoferregung
    \item PQ-Zeit $\approx$ AV-\"Uberleitung
    \item QRS-Komplex $\approx$ Kammererregung
    \item ST-Strecke + T-Welle $\approx$ Erregungsrückbildung
\end{itemize}}

\gCave{Das EKG zeigt nur die \gE{elektrische} Aktivität,
nicht aber die tatsächliche Muskelarbeit an! 

\gEs{Es erlaubt
keine Aussage über die Auswurfleistung! }}

\gFazit{Grundsatz: \gEs{Selbst ein pulsloser Patient kann ein
{\quotedblbase}schönes{\textquotedblleft} EKG haben (allerdings nicht lange)!}}

\subsubsection{EKG -- Ableitungen}
\label{topic:ekg-ableitungen}
Grundsätzlich unterscheidet man zwischen den Extremitätenableitungen
und den Brustwandableitungen. 

\begin{itemize}
    \item Extremitätenableitungen\gZ{: I, II, III, aVL, aVR, aVF}
    \item Brustwandableitungen\gZ{: V1 -- V6}
\end{itemize}

Es haben sich folgende Begriffe eingebürgert:

\begin{description}
    \item[Extremitätenableitung] "`\gT{Rhythmusstreifen}"'
    zur Beurteilung des Herzrhythmus und zum Monitoring
    \item[Extremitätenableitung\,+\,Brustwandableitung]
    "`\gT{12-Kanal-EKG}"' -- Damit können Schäden der
    jeweiligen Herzregion zugeordnet werden. Wenn vorhanden,
    ist das 12-Kanal-EKG (\emph{"`12er"'}) für die
    Diagnostik das Mittel der Wahl. 
\end{description}

\gMarriedNG{Extremitätenableitungen}
{
Die Extremitätenableitungen werden mittels 4 Elektroden abgeleitet. Je
eine Elektrode wird an den Enden der Extremitäten angebracht. Zur
Vereinfachung können die Elektroden auch am \"Ubergang vom Rumpf zu
den Extremitäten angebracht werden (z.B. statt dem rechten Arm an der
rechten Schulter). Wichtig ist, dass die 
Elektrodenpärchen immer auf \gEs{gleicher Höhe} angebracht
werden!

Aus den Informationen der 4 Elektroden erzeugt das
EKG-Gerät \gEs{6 Ableitungen} \gZ{(I, II, III, aVL, aVR
und aVF).}
}{
\begin{itemize}
    \item 4 Elektroden an den Extremitäten
    \item (bzw. am Übergang Rumpf -- Extremitäten)
    \item 6 Ableitungen
\end{itemize}
}





\begin{figure}[H]
\begin{center}
\includegraphics[width=7cm]{./images/geraete/ECGcolor.png}
\gCaptionof{figure}{\gAuthNotesPicLegalOk{}{PD}Extremitätenableitungen
(schwarze Kabel) und Brustwandableitungen (blaue Kabel).
\gSourcePicture{Madhero88 (Wikimedia Commons), Public
domain}\gAuthNotes{GABS: Bessere Beschriftung wäre
wünschenswert, Dringlichkeit: niedrig}}
\end{center}
\end{figure}


\begin{table}[H]
\gCaptionof{table}{Farbe und Positionen der EKG-Elektroden
für die Extremitätenableitung.}

\begin{tabular}{lll}
\gTabHeading{Elektrode} 
& 
\gTabHeading{Position} 
&
\gTabHeading{Alternative Postion}
\\\midrule
\colorbox{red}{\sffamily\bfseries\textcolor{white}{Rot}} & Rechte Hand & Rechte
Schulter
\\
\colorbox{yellow}{\sffamily\bfseries Gelb} & Linke Hand & Linke
Schulter
\\
\colorbox{green}{\sffamily\bfseries Grün} & Linker Fuß  &
Linke Leiste
\\
\colorbox{black}{\sffamily\bfseries\textcolor{white}{Schwarz}} & Rechter Fuß &
Rechte Leiste \\\bottomrule
\end{tabular}
\end{table}

\gMarriedNG{Brustwandableitungen}
{
Die Brustwandableitungen werden mittels 6 Elektroden abgeleitet und
erlauben eine genauere Zuordnung von Störungen zu der jeweiligen
Region des Herzens. Aus den Informationen der 6 Elektroden erzeugt das
EKG-Gerät 6 Ableitungen: \gZ{V1 -- V6}. 
}{
\begin{itemize}
    \item 6 Elektroden am Brustkorb.
    \item 6 Ableitungen
\end{itemize}
}





\begin{table}[H]
\gZ{
\gCaptionof{table}{Position der EKG-Elektroden bei der Brustwandableitung.}

\begin{tabulary}{\linewidth}[H]{CL}
Elektrode
 &
Position
\\\midrule
V1 & 4. Zwischenrippenraum rechts vom Sternum
\\
V2 & 4. Zwischenrippenraum links vom Sternum
\\
V3 & Mitte zwischen V2 und V4
\\
V4 & 5. Zwischenrippenraum links in der Verlängerung der
Schlüsselbein-Mittellinie
\\
V5 & Mitte zwischen V4 und V6
\\
V6 & 6. Zwischenrippenraum in der mittleren Achsellinie
\\\bottomrule
\end{tabulary}
}
\end{table}

\subsection{Pulsoxymetrie}

% \subsubsection{Pulsoxymetrie}
% 
% \index{Pulsoxymetrie}\begin{center}
%  [Warning: Image ignored] % Unhandled or unsupported graphics:
% %\includegraphics[width=3.516cm,height=3.516cm]{AASdev20090228-img13.png}
% \gTempPicMissing{Pulsoxy img13}
% 
% \end{center}
% Das Pulsoxymeter (kurz: Pulsoxy) mißt dden Puls und die
% Sauerstoffsättigung (\index{SpO2}SpO2)des Blutes (Verhältnis
% O2-beladene/ nichtbeladene Erythrozyten). 
% 
% Normalwert 96-99\%, 
% 
% [F078?] Achtung: Ein {\quotedblbase}guter{\textquotedblleft}
% \index{SpO2}Sp\ce{O2}-Wert bedeutet nicht {\quotedblbase}keinen Sauerstoff
% geben{\textquotedblleft}. Die Entscheidung bezüglich Sauerstoffgabe
% berücksichtigt die Diagnose und den Zustand des Patienten, der
% \index{SpO2}Sp\ce{O2}-Wert hat wenig bis gar nichts zu sagen! 
% 
% Achtung: Das Pulsoxymeter verwechselt \ce{O2} und
% \ce{CO}-Moleküle, bei einer \ce{CO}-Vergiftung zeigt das Pulsoxy
% falsch-hohe Sättigungswerte an!
% 
% Weitere Fehlerquellen: kein Signal: Nagellack, Zentralisation
% 
% Das Pulsoxymeter misst die Herzfrequenz und die Sauerstoffsättigung
% (=Sp\ce{O2}) des Blutes, aber es gilt:
% 
% [F078?] Ein {\quotedblbase}guter{\textquotedblleft} Sättigungswert ist
% keine Begründung keinen Sauerstoff zu verabreichen!

Die Pulsoxymetrie wird unter
\prettyref{topic:pulsoxymetrie} besprochen.

\subsubsection{Allgemein}

\gLayout{2}{\gTxBeschreibung}
{

}{}{}{}{}

\gLayout{2}{Funktionen}
{

}{}{}{}{}

\gLayout{2}{Anwendung}
{

}{}{}{}{}

\gLayout{2}{Hygienische Wiederaufbereitung}
{

}{}{}{}{}

\gLayout{2}{\gTxGefahren}
{

}{}{}{}{}

\subsubsection{Gelbes Pulsoxy}

\subsubsection{Rotes Pulsoxy}


\subsection{Blutzuckermessung}

Die Blutzuckermessung wird unter
\prettyref{topic:blutzucker-messung} besprochen.

\subsubsection{AccuCheck}
% TODO : Neufassung
\begin{anfxerror}[author=GABS]{Fehlender Text.}
  Text fehlt! Hier muß neuer Text geschrieben werden!
\end{anfxerror}

\subsection{Fieberthermometer}

% TODO INHALT: Neufassung GER Fieberthermeter
\begin{anfxerror}[author=GABS]{Neufassung GER Fieberthermeter: Fehlender Text.}
  Text fehlt! Hier muß neuer Text geschrieben werden!
\end{anfxerror}

\gLayout{2}{Ohrtermometer}
{~

\gTakeNotesLarge

\subparagraph{Hygienische Maßnahmen}~

\gTakeNotesMedium
}{}{}{}{}

\gLayout{2}{Standardthermometer}
{~

\gTakeNotesLarge

\subparagraph{Hygienische Maßnahmen}~

\gTakeNotesMedium
}{}{}{}{}




\section{Sauerstoff -- \ce{O2}}
\label{topic:sauerstoff}\label{topic:o2}\index{Sauerstoff}\index{O2@\ce{O2}}

\gDefinition{\gT{Sauerstoff} ist ein
lebenswichtiges Gas welches zu 21\% in unserer Atemluft
vorkommt und außerdem ein essentielles
Notfall{\textquotedblright}medikament{\textquotedblleft}.
Das chemische Symbol lautet \gT{\ce{O2}}.}

Sauerstoff darf und muss bei Bedarf vom Sanitäter
verabreicht werden. Die frühzeitige Verabreichung ist wichtig für die Prognose des
Patienten!

% \begin{center}
% \includegraphics[width=4.522cm,height=5.628cm]{./images/rippenspreitzer-02-fett.png}
% \gCaptionof{figure}{NEIN! \gSourcePicture{Rippenspreitzer.de}}
% 
% \end{center}

\subsection{Sauerstoff -- technische Grundlagen}
% TODO : Layout-Anpassung
\begin{anfxwarning}[author=GABS]{Layout-Anpassung notwendig.}
   Dieser Teil entspricht nicht dem Standard-Layout!
\end{anfxwarning}
Sauerstoff wird in Druckflaschen  gelagert. Laut EU-Norm sind diese
Flaschen weiß und mit einem
{\quotedblbase}\gEs{N}{\textquotedblleft} gekennzeichnet.
Das {\quotedblbase}N{\textquotedblleft} steht für {\quotedblbase}neue Kennzeichnung{\textquotedblleft}, da früher die
Flaschen blau gekennzeichnet wurden. \footnote{{\textquotedblright}N{\textquotedblleft} bedeutet in diesem Fall nicht Stickstoff!}

Die Druckflaschen haben unterschiedliche Füllgrößen (Volumina),
gängige Flaschengrößen sind 0,5, 1, 1,5, 2, 5 und
10\,l. Der Sauerstoff wird unter Druck gelagert, wodurch die
Menge des zur Verfügung stehenden Sauerstoffs um ein Vielfaches
gesteigert wird. Der Druck einer gefüllten Flasche beträgt
normalerweise ca. 200\,bar. 

Jede Druckflasche verfügt über ein \gT{Hauptventil} und
ein Standardgewinde. An diesem Gewinde ist entweder ein
weiterführender Druckschlauch oder der \gE{Druckminderer}
einer \gE{Berieselungseinheit} angeschlossen. 

\paragraph{\gT{Berieselungseinheit}}
% TODO : Layout-Anpassung
\begin{anfxwarning}[author=GABS]{Layout-Anpassung notwendig.}
   Dieser Teil entspricht nicht dem Standard-Layout!
\end{anfxwarning}
Die
Berieselungseinheit besteht im Wesentlichen aus einem
\gT{Druckminderer}. Der Druckminderer hat ein
\gE{Abgabeventil}, welches den \gE{Druck drosselt} und eine Abgabe an den Patienten ermöglicht. Je nach Bauart kann der Druckminderer über ein \gE{Regelventil}, ein \gE{Flowmeter}\index{Flowmeter}
und ein \gE{Manometer}\index{Manometer} verfügen. Mit dem
\gT{Regelventil} kann man die Durchfluss- bzw. Abgaberate ein\-stellen, das
\gT{Flowmeter} (Durchflussanzeige) zeigt die aktuelle
Durchflussrate in Liter pro Minute (\gLMin) an. Das
\gT{Manometer} zeigt den Druck in der Sauerstoffflasche an
wenn das Haupt\-ventil geöffnet ist.

\subsection{Heiteres Rechnen mit Sauerstoff}
% TODO : Layout-Anpassung
\begin{anfxwarning}[author=GABS]{Layout-Anpassung notwendig.}
   Dieser Teil entspricht nicht dem Standard-Layout!
\end{anfxwarning}

Mit Sauerstoff lässt sich gut rechnen ...

Die insgesamt verfügbare Menge wird errechnet indem man
die Füllgröße der Flasche (in~l) mit dem
Flaschendruck (in bar) multipliziert.

\begin{equation}
\mbox{Verfügbares Volumen in l} = {\begin{array}{c}
\mbox{Volumen der Flasche in l}\\
\mbox{(Füllgröße)} \end{array}}
\times 
\begin{array}{c}{\mbox{Druck in bar}}\\
\mbox{(Flaschendruck)}
\end{array}
\end{equation}


%oder kurz: 

%\begin{equation}
%l = {Volumen} \times {Druck}
%\end{equation}

Wenn man wissen möchte, wie lange die vorhandene Sauerstoffmenge bei
einer bestimmten Abgaberate reicht, dividiert man einfach
die gesamt verfügbare Menge durch die Abgaberate. 

\gCa{Achtung!} Es ist eine \gEs{ausreichende Reserve}
einzuplanen, in der Flasche sollten min. \gEs{10 bar
Restdruck} verbleiben! 

\gFazit{Es ist eine \gEs{Zeitreserve} und der \gEs{Restdruck} abzuziehen!}

\begin{equation}
\mbox{Zeit} = \frac{\mbox{Verfügbare
Menge}}{\mbox{Abgaberate}}
\end{equation}

%oder kurz: 

%\begin{equation}
%min = \frac{l}{l / min}
%\end{equation}


\minisec{Beispiel}

\begin{quote}
Man hat eine Flasche mit einer Füllgröße von 10\,l. Diese
steht unter einem Druck von 150\,bar: Daraus ergibt
sich, dass in der Flasche 1500\,l \ce{O2} sind.

Der Patient benötigt 5\,\unitfrac{l}{min} \ce{O2} : 

Der Inhalt würde theoretisch für 300 Minuten reichen. Ein
Reservevolumen ist dabei aber noch nicht berücksichtigt!
\end{quote}



\subsection{Umgang mit Sauerstoff und Druckflaschen}

\gLayout{0}{Besondere \gCa{Gefahrenquellen}}
{
Sauerstoff ist \gEs{brandfördernd} und zusammen mit Fett
sogar \gEs{selbstentzündlich}! Neben der Brand- und
Explosionsgefahr birgt der große Druck unter dem die Flaschen
stehen weitere Gefahren: Sollte ein Ventil abbrechen kann die
Flasche eine enorme Kraft entwickeln und zu einer Rakete
werden. Die Kraft einer mit 200\,bar gefüllten Flasche reicht
aus um Mauern zu durchbrechen!

\gCave{
\gCa{Achtung! Vorsicht beim Hantieren -- Zerknall- und
Explosionsgefahr!}
\begin{enumerate}
    \item Sauerstoff ist zusammen mit Fett selbstentzündlich!
    \begin{enumerate}
        \item Die Armaturen und Gewinde sind fettfrei zu
        halten!
        \item Ventile dürfen nicht mit fetten oder öligen Händen
        bedient werden! 
    \end{enumerate}
    \item Das Hantieren mit Feuer und offenem
    Licht ist in der Nähe von Sauerstoffflaschen verboten!
    \item Die Druckflaschen sind immer zu sichern und vor
    Umfallen, etc. zu schützen!
\end{enumerate}
}
}{
\gCa{← Siehe Langtext!}
}{}{}{}

\gLayout{0}{Wechsel und Versorgung nicht benutzter Systeme}
{
Wenn das Sauerstoffsystem nicht benutzt wird muss das \gEs{Hauptventil
geschlossen} und der \gEs{Restdruck im Entnahmeschenkel vollständig abgelassen}
werden. Gleiches gilt für den Wechsel einer Sauerstoffflasche: Es muss sichergestellt sein, dass
\gEs{im Entnahmeschenkel \gCa{kein Überdruck}} herrscht. Vor dem Abmontieren des
Druckminderers oder einer Druckleitung muss dies auf alle Fälle \gEs{extra am
Manometer kontrolliert werden}!

\gCave{\gCa{Lebensgefahr!} Vor der Abmontage eines Druckminderers oder einer
Druckleitung muss die Überdruckfreiheit geprüft werden! } 
}{
\gCa{← Siehe Langtext!}
}{}{}{}

\gFazit{Die Sauerstoffflaschen sind pfleglich zu behandeln!}


\subsection{Verabreichungsarten}
% TODO : Layout-Anpassung
\begin{anfxwarning}[author=GABS]{Layout-Anpassung notwendig.}
   Dieser Teil entspricht nicht dem Standard-Layout!
\end{anfxwarning}

Grundsätzlich unterscheidet man zwischen Berieselung und
Beatmung: 

\subsubsection{Berieselung}
\index{Berieselung}

Die \ce{O2}-Berieselung kann \gEs{nur bei vorhandener
Spontanatmung} durchgeführt werden!

% \marginpar{\gT{Sauerstoff\-brille}}
% \marginpar{\gT{Sauerstoff\-maske}}
% \marginpar{\gT{Sauerstoff\-maske + Reservoir}}

% zwei kurze, in die Nasenöffnungen  eingeführte Kunst\-stoff\-stutzen
% 
% bei ca. 4\,l O2 /min bis 40\% \ce{O2}  in Atemgas  möglich
% 
% 
% Bei behinderter Nasenatmung kann die  Sauerstoffbrille im Mundwinkel
% plaziert werden. 



\begin{table}[H]
\begin{gWideParEnv}
\gCaptionof{table}{Verabreichungsarten von Sauerstoff.}

\begin{tabulary}{\linewidth}[H]{CLLL}
\gTabHeading{Was}
 &
\gTabHeading{Wann}
 &
\gTabHeading{Vorteile}
&
\gTabHeading{Nachteile}
\\\addlinespace\toprule\addlinespace 
\gTabSub{Sauerstoff\-brille}
 &
Spontanatmung

\ce{O2}-Flow {\textless}~5\gLMin
 &
gute Toleranz

fehlendes Engegefühl

Sprechen \& Husten möglich
 &
ungenaue Dosierung

Austrocknung  der Schleimhäute 
\\\addlinespace\midrule\addlinespace
\gTabSub{Sauerstoff\-maske}
 &
Spontanatmung

\ce{O2}-Flow ab 5\gLMin
 &
höhere \ce{O2}-Konzentration
 &
\ce{CO2}-Rückatmung  wenn \ce{O2}-Durchfluß
{\textless}~5\gLMin!) 

für Patienten evt. unangenehm bzw. gewöhnungsbedürftig
\\\addlinespace\midrule\addlinespace
\gTabSub{Sauerstoff\-maske mit Reservoir}
 &
Spontanatmung

\ce{O2}-Flow ab 5\gLMin
 &
Noch höhere \ce{O2}-Konzentration
 &
Wie oben.

Reservoir muss zuerst händisch befüllt werden.
\\\addlinespace\bottomrule
\end{tabulary}
\end{gWideParEnv}
\end{table}

\gDias{Bilderserie: Mittel zur Verabreichung von Sauerstoff.}
{\includegraphics[width=2.5cm]{./images/o2-brille.png}
\gCaptionof{figure}{Sauerstoffbrille}}
{\includegraphics[width=2.5cm]{./images/o2-maske.png}
\gCaptionof{figure}{Sauerstoffmaske}}
{\includegraphics[width=2.5cm]{./images/o2-maske-reservoir.png}
\gCaptionof{figure}{Sauerstoffmaske mit Reservoir}}




\subsubsection{Beatmung}
\index{Beatmung!Grundsätzliches}
% TODO : Layout-Anpassung
\begin{anfxwarning}[author=GABS]{Layout-Anpassung notwendig.}
   Dieser Teil entspricht nicht dem Standard-Layout!
\end{anfxwarning}

Bei \gEs{fehlender} oder unzureichender \gEs{Spontanatmung} (Eigenatmung) wird
eine künstliche Beatmung mittels
\gEs{Beatmungsbeutel} (\gSee{topic:beatmungsbeutel}) oder \gEs{Beatmungsgerät}
durchgeführt. Dabei handelt es sich um eine \gE{Überdruckbeatmung}, d.h. es
wird mit Druck Luft in die Atemwege des Patienten gepumpt. (Im Unterscheid dazu
wird bei der natürlichen Atmung Luft mit Unterdruck in den Patienten
eingesogen.) 

Der notwendige Überdruck kann zu Problemen führen, wie z.B. Überschreiten des
Speiseröhren-Öffnungsdrucks (Magenbeatmung, Magenblähung, Erbrechen) oder
Lungenschäden.



% \gCave{JEDE BEATMUNG SOLLTE MIT RESERVOIRBEUTEL UND 10-15\gLMin
% \ce{O2} DURCHGEF\"UHRT WERDEN!}
% 
% Beatmungsfrequenz beim Erwachsenen: 12-15\,x\,/\,min (Eigenfrequenz)
% 
% Achtung! Zuviel ist schlecht! Nicht zu schnell beatmen ${\rightarrow}$
% Hyperventilationsgefahr!
% 
% \subsection{Sauerstofftherapie}
% Prinzipiell sollte sich die gewählte Therapie an der Diagnose und am
% Schweregrad orientieren!
% 
% Bedenke: 
% 
% \begin{itemize}
% \item Diagnose!
% 
% \item Schockbekämpfung
% 
% \item Ischämie (bei Akutem Koronoarsyndrom, Insult, ..)
% 
% \item Atemstörungen (Fremdkörper, Atemnot, {\dots})
% 
% \end{itemize}

\section{Hilfsmittel für die Beatmung}
\subsection{Beatmungsbeutel}
\label{topic:beatmungsbeutel}
\index{Beatmungs!-beutel}



\gLayout{22}{Material}
{
Eine Beatmung durch Fachpersonal sollte wenn möglich immer mit Hilfmitteln wie
dem \gE{Beatmungsbeutel} oder einem \gE{Beatmungsgerät} durchgeführt werden.
Der Beatmungsbeutel ist das einfachste Hilfmittel welches dem Fachpersonal zur
Verfügung steht und kommt in vielen Bereichen (Rettungsdienst, Krankenhaus,
\ldots) zum Einsatz. Er ermöglichst die kontrollierte oder unterstützende
Beatmung eines Patienten.

Wenn vorhanden soll der Beutel mit einer Sauerstoffzufuhr verbunden werden. Am
einfachsten erfolgt dies über einen Verbindungsschlauch und ein eine
\ce{O2}-Berieselungseinheit, welches mit einem Flow von 15\gLMin{}
eingestellt ist. Mit Hilfe eines Reservoirs, welches Sauerstoff während der
Beatmungspausen zwischenspeichert,  kann die \ce{O2}-Konzentration der
Einatemluft weiter gesteigert werden. 

\gZ{Andere \ce{O2}-Systeme sind am Markt
verfügbar (Demand-Ventile, etc.), aber nur begrenzt verbreitet.}
}{}{./images/geraete/ger-ambu.jpg}{Beatmungsbeutel mit Reservoir, Bakterienfilter und Masken unterschiedlicher Größe.}{}




\begin{table}[H]
\gCaptionof{table}{Mögliche Sauerstoffkonzentrationen bei der Beatmung.}
\begin{tabulary}{\linewidth}{Lr}
\addlinespace 
Mund-zu-Mund Beatmung
&
17\,\%
\\\addlinespace
Beutel-Masken-Beatmung
&
21\,\%
\\\addlinespace
Beutel-Masken-Beatmung mit angeschl. \ce{O2} (15\gLMin)
&
40 - 70\,\%
\\\addlinespace
Beutel-Masken-Beatmung mit angeschl. \ce{O2} (15\gLMin) und
Reservoir &
fast 100\,\%
\\
\end{tabulary}
\hrule
\end{table}


\gLayout{22}{Bestandteile}
{
Ein Beatmungsbeutel besteht im Idealfall aus folgenden Teilen:
\begin{itemize}
	\item Beatmungsmaske 
	\item Bakterienfilter
	\item Ein-/Ausatemventil mit angeschlossenem PEEP-Ventil
	\item Beatmungsbeutel
	\item Reservoir, welches mit einem Verbindungsschlauch
	(\gE{\ce{O2}-Line}) mit einer Berieselungseinheit verbunden
	ist.
\end{itemize}
}{}{./images/geraete/ambu-komplett-zerlegt.jpg}{Bestandteile eines
Beatmungsbeutels.}{}


\subsection{Beatmungsgerät}
\index{Beatmungs!-gerät}
\index{Beatmungs!-gerät!Medumat{\texttrademark}}
\index{Medumat{\texttrademark}}
\index{Medumat{\texttrademark}!Standard a}
\index{Medumat{\texttrademark}!Compact}
\index{Medumat{\texttrademark}!Standard}
\index{Medumat{\texttrademark}!Transport}

\begin{figure}[H]
\begin{center}
\includegraphics[width=7.189cm,height=4.94cm]{./images/beatmungsgeraet-medumat.png}
\gCaptionof{figure}{Beatmungsgerät \emph{Medumat
Standard\textregistered}. \gSourcePicture{ASB
Floridsdorf-Donaustadt}}
\end{center}
\end{figure}




Mit dem Beatmungsgerät \emph{"`Medumat"'} der Fa. Weinmann
kann ein Patient sowohl beatmet als auch berieselt werden
(sofern ein Berieselungsmodul vorhanden ist). 

Dafür gibt es zwei Bedienfelder: Das eine zum Berieseln,
welches vom Sanitäter eingestellt und verwendet werden
kann; das zweite darf nur auf Anweisung eines Notarztes
verwendet werden und wird für das Beatmen benötigt. 

\opt{ORGASBFLO}{
\begin{minipage}{\linewidth}
\ASBFLO{Beim ASB Floridsdorf-Donaustadt kommen folgende Geräte
zum Einsatz:

\begin{itemize}
    \item Hersteller: Fa. Weinmann, Produktfamilie:
    \gT{Medumat} (Notfallbeatmungsgeräte)
    \begin{itemize}
        \item Medumat{\texttrademark} Compact
        \item Medumat{\texttrademark} Standard
        \item Medumat{\texttrademark} Standard a
    \end{itemize}
    \item Hersteller: Fa. Weinmann, Produktfamilie: \gT{Medumat{\texttrademark}}
   
    ((Intensiv-)Transportbe\-atmungs\-ge\-rät\nocite{ThierbachInterhospitaltransfer1})
    \begin{itemize}
        \item Medumat{\texttrademark} Transport
    \end{itemize}
\end{itemize}
}
\end{minipage}
}

\opt{RS}{Test RS}

\begin{figure}[ht]
\includegraphics[width=\linewidth]{./images/geraete/ger-medumat-standard-bedienfeld.jpg}
\gCaptionofDiff{figure}{Bedienfelder Berieselungseinheit \gE{Modul Oxygen} (li.)
und Beatmungsgerät \gE{Medumat{\texttrademark} Standard} (re.).

\gE{Modul Oxygen (li.):} Flowmeter, Ein-Ausschalter und Regelventil. Der
Anschluß links dient der Verbindung der Einheit mit dem Sauerstoffnetz des Fahzeuges.

\gE{Medumat{\texttrademark} Standard (re.):} \gZ{Air-Mix-Schalter,
Ein-Ausschalter, Stellräder für Atemfrequenz, Atemminutenvolumen (!) und Druckbegrenzung,
Beatmungsdruckanzeige, Warnleuchten für Verlegung ("`Stenosis"'), Lösung des
Beatmungsschlauches ("`Disconnection"'), niederen Flaschendruck ("`< 2.7 bar"')
und schwache Batterie.} }
{Bedienfelder Berieselungseinheit \gE{Modul Oxygen}
(li.) und Notfallbeatmungsgerät \gE{Medumat{\texttrademark} Standard} (re.).}
\end{figure}



\subsection{PEEP - Ventil}
\index{PEEP-Ventil}

\gDefinition{PEEP = Positive EndExspiratory Pressure =
positiver end-exspiratorischer Druck}

\gLayout{22}{}
{
Das PEEP-Ventil wird auf der Ausatemöffnung des
Beatmungsbeutels bzw. Beatmungsgerätes aufgesteckt. Es hält
einen positiven Druck in der Lunge am Ende der
Ausatmung aufrecht. Durch Drehen der Kappe kann der
gewünschte PEEP eingestellt werden. Die Einstellung kann an
der angebrachten Skala abgelesen werden. 

Die Anwendung des PEEP ist eine \gEs{ärztliche Maßnahme},
daher muss es \gEs{bei der Anwendung durch einen Sanitäter
immer auf Null stehen}. 

\gZ{Mögliche Einstellungen für das
PEEP-Ventil sind 5 bzw. 10\,cm\,H\textsubscript{2}O (cm
Wassersäule).}

\gCave{Es gibt PEEP-Ventil-Produkte als \gE{Einweg-} oder als
\gE{Mehrweg}material!} 
}
{}{./images/geraete/peep-ventil.jpg}{}{}






\section{Absauggeräte}
\index{Absauggeräte}
\index{Absauggeräte}
\index{Absauggeräte!Accuvac{\texttrademark}}
\index{Absauggeräte!Laerdal Suction Unit{\texttrademark}}
\index{Accuvac{\texttrademark}}
\index{Laerdal Suction Unit{\texttrademark}}

Blut oder Schleim kann vom Fachpersonal mittels Absaugpumpen
und entsprechenden Absaugkatheter naus dem
Mund-Rachen-Raum abgesaugt werden. Es gibt unterschiedliche
Modelle; grundsätzlich wird nach der
Betriebsart unterschieden zwischen \gEs{elektrischer
Absaugeinheit}, \gEs{manueller Absaugpumpe} und als
Sonderfall nur für eine bestimmte Patientengruppe geeignet
der \gEs{Oro-Sauger} oder Schleimabsauger.

\opt{ORGASBFLO}{\ASBFLO{Beim ASB Floridsdorf-Donaustadt werden folgende Geräte
eingesetzt:


\begin{table}[H]
\gCaptionof{table}{ASB921: Eingesetzte Absaugunggeräte.}

\begin{tabulary}{\linewidth}{LLLL}
\gTabHeading{Hersteller}
&
\gTabHeading{Produktfamilie}
&
\gTabHeading{Typ}
&
\gTabHeading{Bemerkungen}
\\\addlinespace\midrule
Weinmann
&
\noindent\gT{Accuvac{\texttrademark}}
&
Rescue
&
rote Deckplatte
\\\addlinespace
&
&
Basic
&
grüne Deckplatte
\\\addlinespace
Laerdal
& 
Laerdal Suction Unit{\texttrademark}
&
&
auslaufend
\\\addlinespace
div. Hersteller
&
Handabsaugpumpe
&
RES-Q-VAC
&
\\\addlinespace
div. Hersteller
&
Orosauger
&
&
\\\addlinespace\bottomrule
\end{tabulary}
\end{table}
}}


\begin{table}[H]
\begin{gWideParEnv}
\gCaptionof{table}{Verschiedene Absaugvorrichtungen.}

\begin{tabular}{p{\linewidth - 10cm}p{4cm}p{4cm}}
\gTabHeading{Elektrische Absaugeinheit}

Älteres Model von der Fa. Laerdal (Laerdal Suction Unit{\texttrademark})

Neueres Model von der Fa. Weimann (Accuvac{\texttrademark})
&
\vspace{0.1pt}
\includegraphics[width=\linewidth]{./images/geraete/lsu-alt.jpg}
&
\vspace{0.1pt}
\includegraphics[width=\linewidth]{./images/geraete/accuvac.jpg}
\\\midrule
\gTabHeading{Manuelle Absaugpumpen}

Handabsaugpumpe

Fußabsaugpumpe
&
\vspace{0.1pt}
\includegraphics[width=\linewidth]{./images/geraete/handabsaugung.jpg}
&
\vspace{0.1pt}
\includegraphics[width=\linewidth]{./images/geraete/fussabsaugung1.jpg}
\\\midrule
\gTabHeading{Oro-Sauger}

Schleimabsauger für Neugeborene
&
\vspace{0.1pt}
\includegraphics[width=\linewidth]{./images/geraete/orosauger.jpg}
\\\midrule
\gTabHeading{Absaugkatheter}

zum Einführen in Mund oder Nase
&
\vspace{0.1pt}
 \includegraphics[width=4cm]{./images/geraete/absaugkatheter.jpg} \\\bottomrule
\end{tabular}
\end{gWideParEnv}
\end{table}

% TODO Erratum: 0.8.0: Absaugung Funktionskontrolle
\gLayout{0}{Funktionskontrolle}
{
Jede Funktionskontrolle umfasst die Überprüfung der\gEs{Vollständigkeit} des
Materials, \gEs{Saugfunktion} bei maximaler und minimaler Leistung und
\gEs{Dichtheit}. Sie muss zumindest bei Dienstbeginn durchgeführt werden.

Bei Geräten, welche über Akkus verfügen, muss die \gEs{Akku-Ladung} und
Saugfunktionsprüfung \gE{ohne Anschluß an eine externe Stromquelle}
erfolgen, d.h. das Gerät muss für die Funktionsprüfung von einer Ladehalterung
abgenommen werden. Die Ladestandsanzeige und evt. vorhandene
\gEs{Warnleuchten} kontrolliert müssen kontrolliert werden. }{
\begin{itemize}
   \item Prüfung von
   \begin{itemize}
      \item Vollständigkeit
      \item Saugleistung
      \item Dichtheit
      \item Ggfs. Akku-Ladung
      \begin{itemize}
         \item Ohne externe Stromquelle!
      \end{itemize}
      \item Warnleuchten
   \end{itemize}
\end{itemize}
}{}{}{}

\gLayout{0}{Absaugkatheter}
{
\index{Absaugkatheter}
Die Absaugkatheter sind steril verpackt und sollten
auch dementsprechend behandelt werden. Wenn 
\gEs{Absaugbereitschaft} hergestellt werden soll, wird nur
der farbige Teil des Katheters ausgepackt und an den
Konnektor bzw.\gT{Fingertip} der Absaugpumpe angeschlossen.
Der restliche Katheter bleibt jedoch noch in der Verpackung. 

Die farbige Kennung gibt
den Durchmesser des Absaugkatheters an. Damit der
Absaugkatheter nicht zu tief eingeführt wird, kann die
einzuführende Katheterlänge durch Abmessen der Distanz
zwischen \gEs{Mundwinkel und Ohrläppchen} bzw. zwischen \gEs{Nasenspitze
und Ohrläppchen} bestimmt werden. Es darf \gEs{maximal auf Sicht} abgesaugt
werden! }{
\begin{itemize}
    \item Auf Sicht.
    \item Mund: Mundwinkel -- Ohrläppchen.
    \item Nase: Mundwinkel -- Nasenspitze
\end{itemize}
}
{}
{}
{}

\gCave{Es darf nur auf Sicht abgesaugt werden!}

\begin{figure}[H]
    \begin{center}
    \includegraphics[width=10cm]{./images/noauth/GER-absaugen.png}
    \gCaptionof{figure}{\gAuthNotesPicLegalNotOk{GER-absaugen.png}{}Absaugung. 
    Es darf nur auf Sicht abgesaugt werden. Die maximale
    Länge ist vom Ohrläppchen bis zur Nasenspitze bzw. zum
    Mundwinkel abzumessen.}
    \end{center}
%%%%%%%%%%%%%%%%%%%%%%%%%%%%%%%%%%%%%%%%%%%%%%%%%%%%%%%%%%%
% TODO Absaugung, Bild:: Überarbeitung: Bild verwirrend
\begin{anfxwarning}[author=KOCH]{Absaugung, Bild:: Überarbeitung: Bild
verwirrend} Absaugung, Bild:: Überarbeitung: Bild verwirrend

KOCH:Bilder verwirrend

Abmessen Nase--Ohr

Absaugen Mund
\end{anfxwarning}
%%%%%%%%%%%%%%%%%%%%%%%%%%%%%%%%%%%%%%%%%%%%%%%%%%%%%%%%%%%
\end{figure}


\section{Transporthilfsmittel}

Das Fachpersonal hat während des gesamten Transports die
Verantwortung für den Patienten, dies gilt auch für gehfähige
Patienten. Sollte der Patient auf Grund seines hohen Alters
oder seiner Kreislaufsituation nicht selbstständig gehen
können, muss er entweder \gEs{sitzend} oder \gEs{liegend}
transportiert werden.

\subsection{Techniken bzw. Geräte für den sitzenden Transport} 

\paragraph{Tragen mit dem Tragering}

Das Tragen mit dem Tragering funktioniert nur bei bewusstseinsklaren und
kooperativen Patienten, außerdem sollte das Stiegenhaus, oder die zurück
zulegende Strecke, breit genug sein. Die Technik eignet sich nur für kurze
Distanzen.

\gLayout{57}{}
{}{}
{
\begin{tabularx}{\linewidth}{XXX}
\includegraphics[width=\linewidth]{./images/motal-aass/tragering1b.jpg}
&
\includegraphics[width=\linewidth]{./images/motal-aass/tragering2b.jpg}
&
\includegraphics[width=\linewidth]{./images/motal-aass/tragering3b.jpg}
\\
Tragering &
Transport mit dem Tragering &
Transport mit dem Tragering 
\\\addlinespace\bottomrule
\end{tabularx}
}{Bilderserie: Transport mit einem Tragering}{}

\gLayout{22}{Umgang mit dem Tragsessel}
{
Der Patient ist immer anzuschnallen.Beim Transport über
Stiegen hat der Patient immer talwärts zu blicken.
}{

}
{./images/noauth/GER-tragsessel.jpg}
{\gAuthNotesPicLegalNotOk{GER-tragsessel.jpg}{}Tragsessel}
{}


\paragraph{Umgang mit einem Rollstuhl}
\index{Rollstuhl}

Siehe Abb. \ref{fig:rollstuhl-umgang}.

~

\subsection{Techniken bzw. Geräte für den liegenden Transport}

\paragraph{Tragetuch}
\index{Tragetuch}

Das Tragetuch (ESE-Tuch) eignet sich ideal für den Transport
von liegenden Patienten wo der Einsatz einer Krankentrage
eventuell aus Platzgründen nicht möglich ist. Es wird auch
gerne zum Umlagern von Patienten oder zum Retten aus einer
Gefahrenzone verwendet. Die Füße des Patienten zeigen
talwärts, in der Ebene wird der Patient wird immer mit den
Füßen in Transportrichtung getragen.

Der Transport mit Hilfe des Tragetuchs ist jedoch nicht
wirbelsäulenschonend. Ist eine Immobilisation der
Wirbelsäule notwendig soll stattdessen eine andere Technik
(Vakuummatratze, etc.) angewandt werden.)

\gLayout{57}{}
{}{}
{
\begin{tabularx}{\linewidth}{XXX}
\includegraphics[width=\linewidth]{./images/motal-aass/tragetuch1b.jpg}
&
\includegraphics[width=\linewidth]{./images/motal-aass/tragetuch2b.jpg}
&
\includegraphics[width=\linewidth]{./images/motal-aass/tragetuch3b.jpg}
\\
Tuch aufbreiten \ldots &
\ldots die näher zum Patient liegende Hälfte
umschlagen\ldots 
& \ldots die obere Hälfte nochmals
umschlagen\ldots 
\\\addlinespace
\includegraphics[width=\linewidth]{./images/motal-aass/tragetuch4b.jpg}
&
\includegraphics[width=\linewidth]{./images/motal-aass/tragetuch5b.jpg}
&
\includegraphics[width=\linewidth]{./images/motal-aass/tragetuch6b.jpg}
\\
\ldots und das Tuch an den Patienten legen. 
&
Patient leicht anheben und Tuch bis zur Hälfte
unterschieben\ldots 
& \ldots Patient schonend zurück legen.
\\\addlinespace
\includegraphics[width=\linewidth]{./images/motal-aass/tragetuch7b.jpg}
&
\includegraphics[width=\linewidth]{./images/motal-aass/tragetuch8b.jpg}
&
\includegraphics[width=\linewidth]{./images/motal-aass/tragetuch9b.jpg}
\\
Von der anderen Seite ebenfalls leicht anheben und
Tuch ausbreiten. 
& 
Zu dritt das Tuch fest anpacken\ldots 
&
\ldots und auf Kommando gemeinsam anheben
\\\addlinespace\bottomrule
\end{tabularx}
}{Bilderserie: Transport mit dem Tragetuch\index{Tragetuch!Bilderserie}}{}

% \begin{figure}[H]
%     \begin{center}
%  TODO BILD: Bilderserie: Tragetuch 
%     \includegraphics[width=10cm]{./images/noauth/GER-tragetuch-tragen.png}
% 
%     \gCaptionof{figure}{\gAuthNotesPicLegalNotOk{GER-tragetuch-tragen.png}{}}
%     \end{center}
% \end{figure}


\subsection{Krankentragen}


\gLayout{22}{Krankentrage mit Fahrgestell}
{\index{Trage}\index{Krankentrage}\index{Fahrtrage}\index{Stollenwerk{\texttrademark}}\index{Ferno!Fahrtrage}
Der Patient ist immer anzuschnallen (Schultern, Becken und Beine!). Beim
Transport über Stiegen hat der Patient immer talwärts zu blicken. 

Wenn ein Patient auf der Trage liegt darf immer nur \gE{ein Ende} gleichzeitig
gehoben oder gesenkt werden, da sonst die Gefahr des Kippens besteht. 

\opt{ORGASBFLO}{Eingesetzte Geräte:

\begin{itemize}
    \item Hersteller: Fa. Stollenwerk
    \item Hersteller: Fa. Ferno
\end{itemize}
}
}{

}
{./images/noauth/GER-fahrtrage-stollenwerk.jpg}
{\gAuthNotesPicLegalNotOk{GER-fahrtrage-stollenwerk.jpg}{}Krankentrage mit Fahrgestell.}
{}

\subsubsection{Krankentrage mit Fahrgestell der Fa. Stollenwerk}
% TODO : Neufassung
\begin{anfxerror}[author=GABS]{Fehlender Text.}
  Text fehlt! Hier muß neuer Text geschrieben werden!
\end{anfxerror}

\begin{compactdesc}
    \item[Hersteller:] Stollenwerk
    \item[Produktfamilie:] Krankentrage, Fahrgestell
    \item[Modell/Typ:] Trage 3002, 3003, 3006; Fahrgestell 2870
    \item[Modellbesonderheiten:] \begin{inparadesc}
    \item[Trage 3002:] Kopfhochstellung. 
    \item[Trage 3003:] Kopf- und Fußhochstellung. 
    \item[Trage 3006:] Kopf- und Fußhochstellung, Bauchdecken-Entlastung. 
    \item[Alle Tragen:] sind mit allen Stollenwerk-Fahrgestellen kompatibel!
\end{inparadesc}   
    \item[Max. Belastbarkeit:] ---\,kg
    \item[Homepage:] \url{http://www.stollenwerk-koeln.de}
    \item[Hygienische Wiederaufbereitung:]
\end{compactdesc}

\gLayout{2}{}
{

}{}{}{}{}

\gLayout{2}{\gTxBeschreibung}
{

}{}{}{}{}

\gLayout{2}{Funktionen}
{

}{}{}{}{}

\gLayout{2}{Anwendung}
{

}{}{}{}{}

\gLayout{2}{Hygienische Wiederaufbereitung}
{

}{}{}{}{}

\gLayout{2}{\gTxGefahren}
{

}{}{}{}{}




\subsubsection{Krankentrage mit Fahrgestell der Fa. Ferno}
% TODO : Neufassung
\begin{anfxerror}[author=GABS]{Fehlender Text.}
  Text fehlt! Hier muß neuer Text geschrieben werden!
\end{anfxerror}

\subsubsection{}



\paragraph{Umlagern von Patienten}\index{Umlagern}

    \begin{center}

    \includegraphics[width=10cm]{./images/noauth/GER-umlagern.png} 

    \gCaptionof{figure}{\gAuthNotesPicLegalNotOk{GER-umlagern.png}{}\label{fig:rollstuhl-umgang}\label{fig:umlagern}Umlagern von Patienten.}
    \end{center}


\section{Geräte zum Retten von Patienten}

\subsection{Schaufeltrage}
\index{Schaufeltrage!Beschreibung}
\label{topic:schaufeltrage-beschreibung}

\gLayout{28}{}
{
Anwendung: \gSee{topic:schaufeltrage-handhabung}

Die Schaufeltrage wird primär für das Umlagern von liegenden
Patienten mit dem Verdacht auf eine Verletzung der
Wirbelsäule verwendet. Das Anwendungsgebiet ist jedoch sehr
vielfältig. Es reicht vom sicheren Transport im unwegsamen
Gelände bis zum Retten aus einer LKW-Kabine.
}{

}
{./images/noauth/GER-schaufeltrage.jpg}
{\gAuthNotesPicLegalNotOk{GER-schaufeltrage.jpg}{}Eine Schaufeltrage}
{}

\subsection{Spineboard}
\label{topic:spineboard-beschreibung}\index{Spineboard!Beschreibung|gZ}

Anwendung: \gSee{topic:spineboard-anwendung}

Das Spineboard ist im Prinzip eine flaches Brett auf dem
der Patient ähnlich dem Schaufeltrage/Vakuummatratze-Duett
gerettet und fixiert werden kann. Diese Technik wurde
ursprünglich vor allem im anglo-amerikanischen Raum
eingesetzt und erfährt auch hierzulande immer größere
Verbreitung.

\ASBFLO{Das Spineboard gehört derzeit nicht zur
Standardausstattung und wird nicht routinemäßig verwendet.}

% \gTakeNotesMedium

\subsection{Rettungskorsett}
\index{Rettungskorsett!Beschreibung}\label{topic:rettungskorsett-beschreibung}

Anwendung: \gSee{topic:rettungskorsett-anwendung}

\gLayout{23}{}
{
Das Rettungskorsett wird primär für das Umlagern von
sitzenden Patienten mit dem Verdacht auf eine Verletzung der
Wirbelsäule verwendet. Es kann aber auch für das Retten
aller sitzenden Patienten aus einer Gefahrenzone verwendet
werden.
}{

}
{./images/geraete/KED.jpg}
{Ein Rettungskorsett.}
{Wikipedia, "`Onthost"', Public domain}




   
\clearpage
\section{Schienungsmaterial}

\subsection{HWS-Schiene}
\label{topic:hws-schiene-beschreibung}\index{HWS-Schiene!Beschreibung}

Anwendung: \gSee{topic:hws-schiene-anlegen}
\gLayout{23}{}
{
Dient zum Ruhigstellen der {\bfseries H}als{\bfseries
W}irbel{\bfseries S}äule (HWS). Es sind verschiedene
Fabrikate im Einsatz (z.B. Stifneck\texttrademark). Es gibt
Schienen mit einer fixen Größe oder moderne Schienen, bei
welchen sich die gewünschte Größe einstellen lässt
(Stifneck select{\texttrademark} bzw. Stifneck Paedi
select{\texttrademark} für Kinder).
}{

}
{./images/noauth/GER-stiffneck.jpg}
{\gAuthNotesPicLegalNotOk{GER-stiffneck.jpg}{}HWS-Schiene.}
{}




\subsection{SAM-Splint}

\label{topic:sam-splint-beschreibung}

Anwendung: \gSee{topic:sam-splint-anwendung}

\gLayout{23}{}
{
Dient am idealsten für das Ruhigstellen von Extremitätenverletzungen mit abnormer Stellung. Da es sich leicht an den
Arm an formen lässt. Es muss jedoch noch mit einer Mullbinde
oder mit Dreiecktücher fixiert werden.
}{

}
{./images/noauth/GER-samsplint.png}
{\gAuthNotesPicLegalNotOk{GER-samsplint.png}{}Sam-Splint.}
{}


\subsection{Vakuumbeinschiene}
\label{topic:vacuumbeinschiene-beschreibung}

\gLayout{23}{}
{
% TODO : Neufassung
\begin{anfxerror}[author=GABS]{Fehlender Text.}
  Text fehlt! Hier muß neuer Text geschrieben werden!
\end{anfxerror}
Dient zum Ruhigstellen von Verletzungen zwischen Knie und Zehen.
}{

}
{./images/geraete/vakuumbeinschiene.jpg}
{\gAuthNotesPicLegalOk{}{}Vakuumbeinschiene.}
{}





\subsection{Vakuummatratze}
\label{topic:vakuummatratze-beschreibung}

\gLayout{23}{}
{
% TODO : Neufassung
\begin{anfxerror}[author=GABS]{Fehlender Text.}
  Text fehlt! Hier muß neuer Text geschrieben werden!
\end{anfxerror}
Die Vakuummatratze wird zum Ruhigstellen von Wirbelsäule, Becken und
Oberschenkel verwendet. }{

}
{./images/noauth/GER-vakuummatratze.png}
{Vakuummatratze.}
{}





%%%
\nocite{RueckerBildatlas1}